\chapter{Compacidad y sucesiones}\label{chap:14}

\section{Por qué buscar una caracterización secuencial}
Las cubiertas abiertas expresan bien el significado de compacidad, pero en problemas concretos solemos disponer de sucesiones: aproximaciones, minimizantes o funciones que queremos hacer converger. En espacios métricos ambas perspectivas son equivalentes. La palabra \emph{métrico} es importante: fuera de esta categoría no se pueden intercambiar sin más compacidad y compacidad secuencial. Usaremos el capítulo \ref{chap:13}, puntos de acumulación y bases numerables. La correspondencia de este capítulo es con la lectura 14 de \cite{valencia}.

\section{Tres formulaciones equivalentes}
\begin{definicion}
Un espacio es \emph{secuencialmente compacto}\index{compacidad!secuencial} si toda sucesión tiene una subsucesión convergente a un punto del espacio. Es \emph{compacto por puntos de acumulación}\index{compacidad!por puntos de acumulacion@por puntos de acumulación} si todo subconjunto infinito tiene un punto de acumulación en el espacio.
\end{definicion}

\begin{lema}\label{ch14:acumulacion}
Si $a$ es punto de acumulación de $A$ en un espacio métrico, cada bola alrededor de $a$ contiene infinitos puntos de $A$ distintos de $a$.
\end{lema}
\begin{proof}
Si $B(a,r)\cap(A\setminus\{a\})=\{b_1,\ldots,b_m\}$ fuera finito, podríamos elegir $0<s<r$ menor que todas las distancias positivas $d(a,b_j)$. Entonces $B(a,s)$ no contendría ningún punto de $A$ distinto de $a$, contradiciendo la acumulación. Si la lista ya es vacía, la contradicción es directa.
\end{proof}

\begin{teorema}\label{ch14:equivalencias}
Para un espacio métrico $X$ son equivalentes:
\begin{enumerate}[label=(\roman*)]
\item $X$ es compacto;
\item $X$ es secuencialmente compacto;
\item todo subconjunto infinito de $X$ tiene un punto de acumulación en $X$.
\end{enumerate}
\end{teorema}
\begin{proof}
\emph{(i)$\Rightarrow$(ii).} Dada $(x_n)$, ponemos $F_n=\cl{\{x_k:k\ge n\}}$. Son cerrados no vacíos y decrecientes. El ejercicio 3 del capítulo \ref{chap:13} proporciona $x\in\bigcap_nF_n$. Elegimos $n_1$ con $d(x_{n_1},x)<1$. Una vez elegido $n_j$, la pertenencia $x\in F_{n_j+1}$ permite escoger $n_{j+1}>n_j$ con $d(x_{n_{j+1}},x)<1/(j+1)$. Esta es una subsucesión que converge a $x$. La barra de cerradura en $F_n$ es indispensable.

\emph{(ii)$\Rightarrow$(iii).} Si $A$ es infinito, elegimos puntos distintos $a_n\in A$. Una subsucesión converge a cierto $a\in X$. Toda bola alrededor de $a$ contiene términos de esa subsucesión distintos de $a$, porque los términos son distintos entre sí y a lo sumo uno puede ser igual a $a$. Así $a$ es punto de acumulación.

\emph{(iii)$\Rightarrow$(ii).} Dada $(x_n)$, si su recorrido es finito, algún valor aparece infinitas veces: una unión finita de conjuntos finitos de índices no podría ser $\N$. Se obtiene una subsucesión constante. Si el recorrido es infinito, tiene un punto de acumulación $x$. Elegido $n_j$, el lema \ref{ch14:acumulacion} permite hallar en $B(x,1/(j+1))$ un valor del recorrido distinto de $x_1,\ldots,x_{n_j}$. Algún índice que lo produzca debe ser mayor que $n_j$. Se construye así una subsucesión convergente.

\emph{(ii)$\Rightarrow$(i).} Primero demostraremos dos hechos. Para cada $r>0$, finitas bolas de radio $r$ cubren $X$. Si esto fallara, elegiríamos $x_1$ y, sucesivamente, $x_{n+1}\notin\bigcup_{j=1}^nB(x_j,r)$. Las distancias entre términos distintos serían al menos $r$, de modo que ninguna subsucesión sería de Cauchy y ninguna convergería.

Segundo, toda cubierta abierta $\mathcal U$ tiene un número de Lebesgue. En caso contrario, para cada $n$ habría un conjunto no vacío $A_n$ de diámetro menor que $1/n$ que no estuviera contenido en ningún miembro de $\mathcal U$. Elegimos $a_n\in A_n$ y una subsucesión $a_{n_j}\to a$. Tomamos $U\in\mathcal U$ y $r>0$ con $B(a,r)\subset U$. Para $j$ grande, $d(a_{n_j},a)<r/2$ y $1/n_j<r/2$. Todo $z\in A_{n_j}$ satisface $d(z,a)<r$, luego $A_{n_j}\subset U$, contradicción.

Sea $\lambda$ un número de Lebesgue. Cubrimos $X$ con finitas bolas de radio $\lambda/3$. Cada una tiene diámetro a lo sumo $2\lambda/3<\lambda$, por lo que está contenida en algún miembro de $\mathcal U$. Estos miembros forman una subcubierta finita. Esta prueba no usa el lema de Lebesgue para compactos para demostrar su propia hipótesis: acabamos de obtener el lema directamente de la compacidad secuencial.
\end{proof}

\begin{corolario}\label{ch14:separable}
Todo compacto métrico es separable y admite una base numerable.
\end{corolario}
\begin{proof}
Para cada $n$ elegimos un conjunto finito $D_n$ cuyos centros de bolas de radio $1/n$ cubran $X$. La unión $D=\bigcup_nD_n$ es a lo sumo numerable. Si $B(x,r)$ es una bola, elegimos $n$ con $1/n<r$ y un centro $a\in D_n$ con $d(x,a)<1/n$. Así $D$ es denso. Las bolas de centro en $D$ y radio racional positivo forman una base: si $B(x,r)$ está dado, elegimos $a\in D$ suficientemente próximo a $x$ y un racional $q$ con $d(a,x)<q<r-d(a,x)$. Entonces $x\in B(a,q)\subset B(x,r)$.
\end{proof}

\section{Ejercicios resueltos}
\begin{ejercicio}[Lectura 14, ejercicio 1]
Si cada $(X_n,d_n)$ es compacto, probar que $\prod_{n\ge1}X_n$ es compacto para la métrica
\[
D(x,y)=\sum_{n=1}^\infty 2^{-n}\frac{d_n(x_n,y_n)}{1+d_n(x_n,y_n)}.
\]
\end{ejercicio}
\begin{solucion}
Si algún factor es vacío, el producto es vacío. En el otro caso, la serie converge porque cada sumando es menor que $2^{-n}$. La función $h(t)=t/(1+t)$ es creciente y subaditiva: para $a,b\ge0$,
\[
\frac{a+b}{1+a+b}=\frac a{1+a+b}+\frac b{1+a+b}
\le\frac a{1+a}+\frac b{1+b}.
\]
Por ello cada $h\circ d_n$ satisface la desigualdad triangular y también la suma $D$. La positividad, simetría y separación de puntos se comprueban término a término.

Sea $(x^{(k)})$ una sucesión del producto. Extraemos una subsucesión para la cual converge la primera coordenada; de esta, otra para la segunda; continuamos inductivamente. Escribamos $k_j^{(n)}$ los índices de la subsucesión elegida en el paso $n$, ordenados crecientemente, con cada lista contenida en la anterior. Los índices diagonales $k_n^{(n)}$ son estrictamente crecientes: el $(n+1)$-ésimo índice de la lista siguiente es al menos el $(n+1)$-ésimo de la anterior, y este supera su $n$-ésimo. Para cada coordenada fija $i$, la cola de la diagonal pertenece a la subsucesión del paso $i$ y converge a un punto $x_i\in X_i$.

Sea $x=(x_i)$. Dado $\eps>0$, elegimos $N$ con $\sum_{i>N}2^{-i}<\eps/2$. En las primeras $N$ coordenadas la convergencia permite hacer que la suma finita de distancias transformadas sea menor que $\eps/2$. La cola ya está controlada independientemente del índice. Así $D(x^{(k_n^{(n)})},x)<\eps$ para $n$ grande. El producto es secuencialmente compacto y el teorema \ref{ch14:equivalencias} concluye la prueba. Esta estimación prueba también que la convergencia en $D$ equivale a la convergencia en cada coordenada: la otra dirección se obtiene acotando cada sumando por $D$.
\end{solucion}

\begin{ejercicio}[Lectura 14, ejercicio 2]
Sea $X$ compacto y $U\subset X\times X$ abierto que contiene la diagonal $\Delta_X$. Probar que existe $r>0$ tal que $d(x,y)<r$ implica $(x,y)\in U$.
\end{ejercicio}
\begin{solucion}
Si no existiera tal $r$, elegiríamos $x_n,y_n$ con $d(x_n,y_n)<1/n$ y $(x_n,y_n)\notin U$. Una subsucesión $x_{n_j}$ converge a $x$; la desigualdad $d(y_{n_j},x)\le1/n_j+d(x_{n_j},x)$ muestra que $y_{n_j}\to x$. Entonces $(x_{n_j},y_{n_j})\to(x,x)\in U$. Como $U$ es abierto, la sucesión pertenece eventualmente a $U$, contradicción. El argumento convierte cercanía a la diagonal en una cota uniforme.
\end{solucion}

\begin{ejercicio}[Lectura 14, ejercicio 3]
Probar que todo espacio métrico no compacto contiene un subconjunto infinito numerable, cerrado y discreto.
\end{ejercicio}
\begin{solucion}
Existe una sucesión sin subsucesiones convergentes. Ningún valor aparece infinitas veces; su recorrido es infinito. Elegimos inductivamente una subsucesión de valores distintos y llamamos $A$ a su recorrido. Si $A$ tuviera un punto de acumulación, el lema \ref{ch14:acumulacion} permitiría extraer, respetando el orden de los índices, una subsucesión convergente, contradicción. Por tanto $A$ no tiene puntos de acumulación. Todo punto de $\cl A\setminus A$ sería uno de acumulación; así $A$ es cerrado. Para cada $a\in A$, que $a$ no sea de acumulación significa que existe $r_a>0$ con $B(a,r_a)\cap A=\{a\}$. Esta es precisamente la discreción de la topología relativa.
\end{solucion}

\begin{ejercicio}[Lectura 14, ejercicio 4]
Probar que $X$ es compacto si y solo si toda función continua $X\to\R$ es acotada.
\end{ejercicio}
\begin{solucion}
La dirección directa se sigue de que las imágenes compactas en $\R$ son acotadas. Para la inversa, si $X$ no es compacto, el ejercicio anterior proporciona $A=\{a_n:n\ge1\}$ cerrado y discreto. La función $a_n\mapsto n$ es continua en $A$ y el teorema de Tietze, en su versión real no necesariamente acotada, permite extenderla. Damos también una construcción explícita para no ocultar esa versión del teorema.

Elegimos
\[
0<r_n<\min\{1/n,\tfrac13\dist(a_n,A\setminus\{a_n\})\},\qquad
\varphi_n(x)=\max\{0,1-d(x,a_n)/r_n\}.
\]
La distancia del centro al resto es positiva por discreción. La familia de soportes $\supp\varphi_n\subset\overline B(a_n,r_n)$ es localmente finita. En efecto, si $x\notin A$, $\delta=\dist(x,A)>0$; para $n$ grande, $r_n<\delta/3$, y esas bolas no tocan $B(x,\delta/3)$. Si $x=a_k$, poniendo $\delta=\dist(a_k,A\setminus\{a_k\})>0$, toda bola con $n\ne k$ queda a distancia mayor que $2\delta/3$ de $a_k$, ya que $r_n<d(a_n,a_k)/3$. Por tanto
\[
f(x)=\sum_{n=1}^\infty n\varphi_n(x)
\]
es localmente una suma finita de funciones continuas, luego es continua. Además $f(a_n)=n$. Es una función continua no acotada, contradiciendo la hipótesis.
\end{solucion}

\begin{ejercicio}[Lectura 14, ejercicio 5]
Sea $1\le p<\infty$. Probar que un subconjunto cerrado y acotado $A\subset\ell^p$ es compacto si y solo si
\[
\forall\eps>0\ \exists N\ \forall x\in A:\quad
\sum_{k=N}^\infty |x_k|^p<\eps.
\]
Deducir que $K=\{x\in\ell^2:0\le x_n\le1/n\}$ es compacto.
\end{ejercicio}
\begin{solucion}
\emph{Necesidad.} Sea $A$ compacto. Fijamos $\eps>0$ y elegimos $\delta>0$ con $2^p\delta^p<\eps$. Finitas bolas de radio $\delta$, centradas en $a^{(1)},\ldots,a^{(m)}\in A$, cubren $A$. Cada centro pertenece a $\ell^p$, por lo que para un $N$ común a todos ellos sus colas tienen suma menor que $\delta^p$. Dado $x\in A$, elegimos un centro a distancia menor que $\delta$ y usamos $|u+v|^p\le2^{p-1}(|u|^p+|v|^p)$:
\[
\sum_{k=N}^\infty|x_k|^p
\le2^{p-1}\left(\norm{x-a^{(j)}}_p^p+
\sum_{k=N}^\infty|a_k^{(j)}|^p\right)<2^p\delta^p<\eps.
\]

\emph{Suficiencia.} Sea $(x^{(j)})$ una sucesión en $A$ y $M$ una cota para sus normas. Como cada coordenada está acotada por $M$, el procedimiento diagonal del ejercicio 1, usando Bolzano--Weierstrass en $\R$ o $\C$, produce una subsucesión que converge coordenada a coordenada a $x=(x_k)$. Para todo $m$,
\[
\sum_{k=1}^m|x_k|^p=\lim_j\sum_{k=1}^m|x_k^{(j)}|^p\le M^p.
\]
Las sumas parciales son crecientes y acotadas, por lo que $x\in\ell^p$. Dado $\eta>0$, elegimos $N$ de modo que las colas de todos los elementos de $A$ sean menores que $\eta$. Pasando al límite en cada suma finita de la cola y luego tomando el supremo de esas sumas obtenemos $\sum_{k=N}^\infty|x_k|^p\le\eta$. Así
\[
\norm{x^{(j)}-x}_p^p
\le\sum_{k<N}|x_k^{(j)}-x_k|^p+2^p\eta.
\]
Primero hacemos $\eta$ tan pequeño como necesitemos y después hacemos pequeña la suma finita mediante la convergencia coordenada. Se obtiene convergencia en norma. Puesto que $A$ es cerrado, $x\in A$, y la compacidad secuencial prueba la compacidad.

Para $K$, cada condición $0\le x_n\le1/n$ es cerrada porque la evaluación de coordenada es continua. Además $\norm{x}_2^2\le\sum_n1/n^2<\infty$, y las colas están uniformemente acotadas por $\sum_{k\ge N}1/k^2\to0$. El criterio se aplica. La condición de colas es la parte de la compacidad que la mera acotación en dimensión infinita no detecta.\index{cubo de Hilbert}
\end{solucion}

\begin{ejercicio}[Lectura 14, ejercicio 6]
Sea $W$ una vecindad compacta de $x$ y $f:W\to[0,\infty)$ continua, con $f(x)=0$ y $f(y)>0$ para $y\in W\setminus\{x\}$. Probar que para toda vecindad $V$ de $x$ existe $r>0$ tal que $\{y\in W:f(y)<r\}\subset V$.
\end{ejercicio}
\begin{solucion}
El subnivel debe tomarse en $W$, que es el dominio de $f$; escribir $y\in X$ sin extender $f$ no define un conjunto. Elegimos un abierto $O$ con $x\in O\subset V$. El conjunto $W\setminus O$ es compacto. Si es vacío, cualquier $r>0$ funciona. Si no, $f$ alcanza en él un mínimo $m$. Ninguno de sus puntos es $x$, luego $m>0$. Tomamos $r=m$. Un punto de $W$ con $f(y)<r$ no puede estar en $W\setminus O$, y por ello pertenece a $O\subset V$. El argumento explica por qué la positividad punto a punto se vuelve una cota positiva uniforme lejos del único cero.
\end{solucion}

\section{Idea que conviene conservar}
En un compacto métrico siempre podemos extraer una subsucesión convergente. Para reconocer compactos de sucesiones o funciones conviene buscar dos ingredientes: control de finitas coordenadas y una estimación uniforme de la cola. El primero por sí solo no garantiza convergencia en la norma del espacio.
